\chapter{The WCCI Grid and the Design Carried Onto It}
\label{ch:wcci-transfer}

The controlled method comparisons so far were developed on \texttt{bus14}.
That grid is small enough to iterate on but too small to claim anything about
scale: fourteen substations, twenty lines, and three agents. The preliminary
WCCI stress test of Section~\ref{sec:sparse-wcci-stress-test} reproduced several
sparse-control mechanisms with an MLP actor on the larger grid, but even the
matched baseline was weak. Because that pilot changed grid size, action space,
and maintenance at once, it could not identify the source of failure.
\\
\\
This chapter turns that pilot into a controlled transfer study. It moves to the
L2RPN WCCI 2020 grid, more than twice as many substations and three times as many
lines, and asks the question that motivates a graph encoder in the first place.
If the actor reads the grid as a graph rather than as a flat vector, its weights
are independent of how many substations exist, and a representation learned on
\texttt{bus14} should carry over.
\\
\\
Before interpreting that transfer, the two environments have to be made
comparable and the encoder's inputs have to be checked. This chapter does both.
The comparison turns out to matter more than expected: the graph features the
encoder consumes are raw physical quantities, and they do not follow the same
distribution on the two grids.

\section{The target environment}
\label{sec:wcci-environment}

Table~\ref{tab:wcci-env-comparison} contrasts the source and target grids. The
agent partitions follow the same principle in both cases, contiguous electrical
areas, but WCCI is split into four regions rather than three, and the regions
are markedly unequal: one agent controls seventeen substations while another
controls five.

\begin{table}[H]
  \centering
  \small
  \caption{Source and target environments. Both use topology actions,
  decentralized agents and an 80/20 train/test chronic split with split seed 0.}
  \label{tab:wcci-env-comparison}
  \setlength{\tabcolsep}{6pt}
  \begin{tabular}{L{0.30\textwidth}L{0.28\textwidth}L{0.30\textwidth}}
    \toprule
    & \texttt{bus14} & \texttt{bus36\_wcci\_nomaint} \\
    \midrule
    Grid2Op environment & \texttt{l2rpn\_case14\_sandbox} & \texttt{l2rpn\_wcci\_2020} \\
    Substations & 14 & 36 \\
    Lines & 20 & 59 \\
    Agents & 3 & 4 \\
    Chronics & 1{,}004 & 2{,}880 \\
    Episode length & up to 8{,}064 steps & up to 8{,}062 steps \\
    Maintenance & none & disabled (Section~\ref{sec:wcci-nomaint}) \\
    \bottomrule
  \end{tabular}
\end{table}


\subsection{Removing maintenance}
\label{sec:wcci-nomaint}

Of the environments Grid2Op distributes, \texttt{l2rpn\_case14\_sandbox} is the
only one above the toy scale that ships without scheduled maintenance. Every
larger environment, WCCI 2020 included, injects planned line outages. That is a
problem for a transfer study: a policy trained without maintenance and evaluated
with it faces a hazard it has never seen, and the two observation spaces differ
as well, since the maintenance countdown features are only present when
maintenance exists.
\\
\\
The two settings are therefore aligned by removing maintenance from WCCI rather
than inventing a schedule for \texttt{bus14}, for which Grid2Op provides no
specification.
\\
\\
The effect is measured on the complete held-out test split using the same 576
chronics and a do-nothing policy in both settings.

\begin{table}[H]
  \centering
  \small
  \caption{Do-nothing performance on the full WCCI test split. The chronics and
  policy are identical; only scheduled maintenance differs.}
  \label{tab:wcci-maintenance-effect}
  \setlength{\tabcolsep}{7pt}
  \begin{tabular}{lrrr}
    \toprule
    Setting & Mean survival & Median survival & Complete episodes \\
    \midrule
    Maintenance enabled  & 27.3\% & 12.1\%  & 39/576 (6.8\%) \\
    Maintenance disabled & 57.3\% & 100.0\% & 303/576 (52.6\%) \\
    \midrule
    Difference            & +30.0 pp & +87.9 pp & +264 (+45.8 pp) \\
    \bottomrule
  \end{tabular}
\end{table}



\subsection{Action space}
\label{sec:wcci-action-space}

The unreduced WCCI action space is not trainable. Enumerating topology actions over each agent's region gives the sizes in
the first column of Table~\ref{tab:wcci-action-space}: nearly 67{,}000 actions in
total, of which agent~1 alone accounts for 65{,}642. All experiments therefore
use the reduced action space obtained by the brute-force outcome procedure, at
$k=256$ per agent.

\begin{table}[H]
  \centering
  \small
  \caption{Per-agent action-space sizes on \texttt{bus36\_wcci\_nomaint}. Agents
  0 and 2 are already exhaustive at $k=256$, so the reduction only binds on
  agents 1 and 3.}
  \label{tab:wcci-action-space}
  \setlength{\tabcolsep}{8pt}
  \begin{tabular}{lrrr}
    \toprule
    Agent & Substations & Full & Reduced ($k=256$) \\
    \midrule
    \texttt{agent\_0} & 5  & 77     & 77  \\
    \texttt{agent\_1} & 5  & 65{,}642 & 256 \\
    \texttt{agent\_2} & 9  & 127    & 127 \\
    \texttt{agent\_3} & 17 & 1{,}119  & 256 \\
    \midrule
    Total             & 36 & 66{,}965 & 716 \\
    \bottomrule
  \end{tabular}
\end{table}

\noindent The choice of $k$ was made by evaluating a one-step simulator-greedy policy over
each candidate reduction on the same fifty chronics
(Table~\ref{tab:wcci-greedy-sizes}). Performance saturates at $k=256$: the
difference between $k=256$ and $k=1024$ is four chronics won against three lost,
which is noise, while both clearly beat $k \le 128$. Since $k=256$ is the
smallest reduction in the saturated group, it gives the smallest action heads
for no measurable loss in what the action set can achieve.

\begin{table}[H]
  \centering
  \small
  \caption{One-step simulator-greedy policy over each reduced action space, on
  the same fifty held-out chronics of \texttt{bus36\_wcci\_nomaint}. The
  do-nothing row is the common reference; no reduction was evaluated against a
  different chronic set.}
  \label{tab:wcci-greedy-sizes}
  \setlength{\tabcolsep}{8pt}
  \begin{tabular}{lrrr}
    \toprule
    Policy & Mean survival & Full-episode survival & Episodes worse than do-nothing \\
    \midrule
    Do nothing        & 0.560 & 52\% & --- \\
    Greedy, $k=32$    & 0.679 & 60\% & 5 \\
    Greedy, $k=64$    & 0.819 & 64\% & 3 \\
    Greedy, $k=128$   & 0.847 & 76\% & 3 \\
    Greedy, $k=256$   & \textbf{0.976} & \textbf{92\%} & 0 \\
    Greedy, $k=512$   & 0.883 & 82\% & 2 \\
    Greedy, $k=1024$  & 0.953 & 90\% & 0 \\
    \bottomrule
  \end{tabular}
\end{table}

\noindent These two rows bracket the transfer experiments. A learned policy that lands
near 0.56 mean survival has learned nothing beyond doing nothing; one
approaching 0.9 is doing real work. The greedy row is not a strict upper bound,
since it queries the simulator at every decision step and is myopic, but it is a
strong reference.


\section{Corrected inputs}
\label{sec:wcci-corrected-inputs}

Chapter~\ref{ch:encoder-inputs} measures the encoder's inputs on both grids and
settles two corrections, both applied in every run below: the power channels are
divided by each grid's own maximum generator rating
(Section~\ref{sec:wcci-physical-scaling}), and the absolute voltage angle is
replaced by the drop along each line
(Section~\ref{sec:wcci-angle-representation}). Running standardization is left
off. Together these bring the angle channel into the band occupied by
\texttt{rho} and the power spreads to within 15\,\% of parity
(Table~\ref{tab:wcci-corrected}).


\section{Which screened design to carry forward}
\label{sec:wcci-architecture-choice}

Chapter~\ref{ch:experiments} settled five design questions on \texttt{bus14},
each producing a configuration worth keeping. Screens~A and~B fixed the encoder
depth, width and pooling, giving the plain busbar graph
\texttt{bus\_e0n0v0}. Screen~C augmented it, and the augmentation that paid was
an edge between busbars of the same substation, \texttt{bus\_e1n0v0}. Screen~D
moved to the heterogeneous representation and screened message direction, where
\texttt{het\_ga2b\_lb2a} and the baseline \texttt{het\_gbi\_lbi} came out
ahead. Screen~E was a different line of attack --- heterogeneous with line nodes,
and an action head that scores each candidate from the nodes it touches instead
of from a fixed-width layer over an action index --- and it worked well enough on
\texttt{bus14} to keep.
\\
\\
Which of them should the transfer study use? The question is answerable directly,
without any appeal to transfer: re-instantiate each on
\texttt{bus36\_wcci\_nomaint} and train it there from scratch.
Table~\ref{tab:wcci-proxy-test} does that under one protocol --- $k=64$, 15M
steps, random initialisation, no transferred weights anywhere.

\begin{table}[H]
  \centering
  \small
  \caption{Each screen's surviving configuration, retrained from scratch on WCCI
  at $k=64$, against the \texttt{bus14} survival of the identical model.
  \texttt{het\_gbi\_lbi} is the one cell of Screen~D's $3\times3$ factorial
  without a rerun under the corrected graph construction, so its \texttt{bus14}
  entry is empty rather than filled from the pre-correction screen. One seed per
  cell; Screens~A--D were also run at $k=256$ with the same outcome.}
  \label{tab:wcci-proxy-test}
  \setlength{\tabcolsep}{5pt}
  \input{tables/wcci_proxy_test}
\end{table}

\textbf{Screen~E is the only one that produces a usable policy.} It clears the
do-nothing floor on the difficult cohort in $5$ of its $8$ cells, with a median
of $9.61\,\%$, a best of $18.46\,\%$ and two chronics carried to the end.
Screens~A--D clear it in $0$ of $4$, with a best of $6.80\,\%$ and no rescues at
all --- across every evaluation of those four configurations, including their
best-test checkpoints, not one of the 24 difficult chronics is completed. They do
not merely underperform: they break what do-nothing preserves, keeping a median
of $7.5$ of the 26 easy chronics.
\\
\\
\textbf{\texttt{bus14} score does not explain the gap.} Screens~A--D span
$95.45$--$99.12\,\%$ there and Screen~E spans $93.40$--$100.00\,\%$: overlapping
ranges, opposite outcomes on the target grid. The ranking is informative only
\emph{within} a design --- inside Screen~E, \texttt{bus14} rank orders the WCCI
result at Spearman $+0.74$ ($p=0.037$), and its best \texttt{bus14} cell is also
its best on WCCI --- and carries nothing \emph{between} designs.
\\
\\
\textbf{Screen~E is also the only one that can be transferred whole.}
Screens~A--D use an action head with one output per action index, a shape fixed
by a single grid's action list; their encoders can cross a grid boundary but
their actors cannot. Screen~E's head applies shared weights to a variable-length
candidate set, so its parameters do not depend on the action count and
\texttt{agent\_0}'s scorer can initialise every agent on any grid.
\\
\\
Those two properties are independent --- one is measured in
Table~\ref{tab:wcci-proxy-test}, the other follows from parameter shapes --- and
they point the same way. The remainder of this chapter therefore studies
Screen~E alone, and asks how far its portability can actually be exploited.
\\
\\
Two gaps are worth recording. \texttt{het\_gbi\_lbi} has no corrected
\texttt{bus14} number and is the best of the four Screens~A--D configurations on
WCCI, so the one cell that cannot be anchored is the one carrying most weight
above. And the two classes differ in graph schema as well as action head, because
\texttt{common/action\_metadata.py} requires
\texttt{gnn\_graph\_type = "heterogeneous\_line"} for candidate scoring; the
comparison is therefore between configurations, not between action heads in
isolation.


\begin{tcolorbox}[
  colframe=red!75!black,
  colback=red!5!white,
  title=Chapter 8 Summary: The WCCI Grid and the Design Carried Onto It]
\textbf{Purpose:} To define the target environment precisely enough that a
transfer result is attributable, and to decide --- on the target grid itself,
without appealing to transfer --- which screened design is worth studying there.
\\
\\
\textbf{Content:} WCCI is $2.6\times$ \texttt{bus14} in substations and
$3\times$ in lines, over four unequal agent regions. Maintenance is removed
rather than invented for \texttt{bus14}, and the action space is reduced from
nearly $67{,}000$ topology actions ($65{,}642$ on one agent alone) to $k=256$
per agent. Every configuration surviving Screens~A--E is then retrained directly
on WCCI from random initialisation, $k=64$, 15M steps, no transferred weights.
\\
\\
\textbf{Findings:}
\begin{itemize}
  \item Only Screen~E produces a usable policy: above the do-nothing floor on
  the difficult chronics in $5$ of $8$ cells, median $9.61\,\%$, best
  $18.46\,\%$. Screens~A--D manage $0$ of $4$, best $6.80\,\%$, and never carry
  a chronic to the end.
  \item A \texttt{bus14} score does not explain the gap: the two classes occupy
  overlapping ranges ($95.45$--$99.12\,\%$ against $93.40$--$100.00\,\%$) and
  land on opposite sides of the floor. The ranking orders results only
  \emph{within} a design (Spearman $+0.74$ inside Screen~E).
  \item Screen~E is also the only design transferable whole: the others index a
  fixed action list, so their encoders can cross a grid boundary but their
  actors cannot.
\end{itemize}
\textbf{Takeaway:} The best learner on the target grid and the only portable
design are the same configuration, and the two facts are independent --- one
measured, one following from parameter shapes. Chapter~\ref{ch:screen-e-transfer}
therefore studies Screen~E alone. Two gaps stand: \texttt{het\_gbi\_lbi} has no
corrected \texttt{bus14} number despite being the best of Screens~A--D on WCCI,
and the graph schema moves with the action head, so the comparison is between
configurations rather than action heads.
\end{tcolorbox}
